\subsection{Stable extensions approaching the split comparison bundle}
\label{subsec:hirzebruch-stable-extensions}
\index{Hirzebruch surface!stable extensions}

The bundles near $E_0$ that we construct will also be slope-stable for a fixed polarization. The split bundle $E_0$ is unstable, and the ample bundle $E_{\mathrm{amp}}$ need not be stable. Nonsplit extensions of $L_2$ by $L_1$ have the same Chern classes as $E_0$, are stable, and converge to $E_0$ as the extension class tends to zero; they are not ample, however. The next subsection combines them with $E_{\mathrm{amp}}$.

Keep the Hirzebruch surface $X=\mathbb F_e$, its ruling and the line bundles of \cref{subsec:hirzebruch-ample-comparison}. Thus
\[
 L_1=\mathcal O_X(5S-f),\qquad
 L_2=\mathcal O_X(-S+(12e+5)f),\qquad E_0=L_1\oplus L_2.
\]
We use the very ample polarization $P=4H-f=4S+(12e+3)f$. For a torsion-free coherent sheaf $E$ of positive rank on $X$, put
\[
 \deg_PE=c_1(E)\cdot P,
 \qquad \mu_P(E)=\frac{\deg_PE}{\operatorname{rank}E}.
\]
A vector bundle $V$ is \emph{$\mu_P$-stable} if $\mu_P(M)<\mu_P(V)$ for every coherent subsheaf $M\subset V$ with $0<\operatorname{rank}M<\operatorname{rank}V$. Since $\mu_P(L_2)>\mu_P(E_0)$ by the computation below, the split bundle $E_0$ is not $\mu_P$-stable.

\Needspace{5\baselineskip}
\begin{proposition}\label{prop:hirzebruch-stable-extensions}
\index{Vector bundle!slope-stable extensions}
The extension space has dimension
\[
 \dim\operatorname{Ext}^1_X(L_2,L_1)=105e+35.
\]
Every nonsplit extension
\[
 0\longrightarrow L_1\longrightarrow V\xrightarrow{q}L_2
 \longrightarrow0
\]
is a locally free, $\mu_P$-stable bundle with $c_1(V)=4H$ and $c_2(V)=13H^2$.
\end{proposition}
\begin{proof}
Extensions of line bundles are locally free: the quotient is free on a small trivializing neighborhood, so the sequence splits there. Their classes form
\[
 \operatorname{Ext}^1_X(L_2,L_1)
 =H^1(X,\mathcal O_X(6S-(12e+6)f)).
\]
By \cref{lem:hirzebruch-ruling-sections}, the higher direct image of this line bundle vanishes and its direct image is the sum of the seven line bundles $\mathcal O_{\mathbb P^1}(-12e-6-ie)$, $0\leq i\leq6$. Their first cohomology groups have total dimension
\[
 \sum_{i=0}^6(12e+5+ie)=105e+35>0.
\]
The Chern classes follow from the extension sequence. Direct computations of intersection numbers give
\[
 \mu_P(L_1)=40e+11,\qquad \mu_P(V)=40e+14,
 \qquad \mu_P(L_2)=40e+17.
\]

It suffices to test line-bundle subsheaves $M\hookrightarrow V$. Indeed, a rank-one torsion-free subsheaf on a smooth surface has a line bundle as its double dual, and its inclusion extends across the finitely many missing points by Hartogs extension in local bundle frames. The extended map remains injective and has the same first Chern class. This argument does not require the quotient to be locally free.

If the composite $M\to L_2$ is zero, then $M$ maps nontrivially into $L_1$, so $\deg_P M\leq\mu_P(L_1)<\mu_P(V)$. Otherwise its zero divisor $D$ is effective and $M=L_2(-D)$. If $D=0$, the composite is an isomorphism and its inverse splits $q$, contrary to the hypothesis. Hence $D\neq0$, and
\[
 \deg_P M=\mu_P(V)+3-D\cdot P.
\]
Every effective class is $aS+bf$ with $a,b\geq0$, and
\[
 S\cdot P=8e+3,\qquad f\cdot P=4.
\]
For $e\geq1$, every nonzero effective divisor has $P$-degree at least four, which proves the strict slope inequality.

For $e=0$, the only nonzero effective class of $P$-degree at most three is $S$. A lift of multiplication by a nonzero section $s\in H^0(X,\mathcal O_X(S))$ from $L_2(-S)$ to $V$ would make the pullback of the nonzero extension class vanish. That pullback is the image of the class under the map
\[
 H^1(X,\mathcal O_X(6S-6f))
 \xrightarrow{\ \cdot s\ }H^1(X,\mathcal O_X(7S-6f)).
\]
On $X=\mathbb P^1\times\mathbb P^1$, order the factors so that $\mathcal O_X(S)=\mathcal O(1,0)$ and $\mathcal O_X(f)=\mathcal O(0,1)$. The K\"unneth formula identifies this map with
\[
 \bigl[H^0(\mathbb P^1,\mathcal O(6))
 \xrightarrow{\ \cdot s\ }H^0(\mathbb P^1,\mathcal O(7))\bigr]
 \otimes\operatorname{Id}_{H^1(\mathbb P^1,\mathcal O(-6))}.
\]
Multiplication by a nonzero section is injective on global sections, and tensoring with a vector space preserves injectivity. Thus the nonzero extension class cannot pull back to zero. This excludes the possible equality of slopes and proves stability also for $e=0$.
\end{proof}

\Needspace{5\baselineskip}
\begin{proposition}\label{prop:hirzebruch-extension-prioritary}
\index{Prioritary bundle!extensions}
\index{Vector bundle!negative fibre quotient}
Every extension of $L_2$ by $L_1$, including the split extension, is prioritary for the ruling and is not ample. Its restriction to every ruling fiber is $\mathcal O_{\mathbb P^1}(5)\oplus\mathcal O_{\mathbb P^1}(-1)$.
\end{proposition}
\begin{proof}
Put $D=K_X+f=-2S-(e+1)f$. The four line bundles $L_i^*\otimes L_j(D)$ have classes
\[
 D,\quad D,\quad 4S-(13e+7)f,\quad -8S+(11e+5)f.
\]
Each has a negative coefficient in the basis $S,f$, and hence no nonzero section. Applying $\operatorname{Hom}_X(L_i,\cdot)$ to the extension twisted by $D$ gives $\operatorname{Hom}_X(L_i,V(D))=0$ for $i=1,2$. Applying $\operatorname{Hom}_X(\cdot,V(D))$ to the original extension then gives $\operatorname{Hom}_X(V,V(D))=0$. By Serre duality this is the prioritary condition $\operatorname{Ext}^2_X(V,V(-f))=0$.

Restriction to a ruling fiber is exact, because the quotient is locally free. The restricted line bundles have degrees five and minus one, and the restricted extension splits, since
$\operatorname{Ext}^1_{\mathbb P^1}(\mathcal O(-1),\mathcal O(5))
=H^1(\mathbb P^1,\mathcal O(6))=0$.
The quotient of negative degree on this curve excludes ampleness.
\end{proof}

\Needspace{5\baselineskip}
\begin{proposition}\label{prop:hirzebruch-explicit-stable-degeneration}
\index{Dolbeault operator!smooth convergence}
\index{Vector bundle!stable degeneration}
On the smooth bundle underlying $E_0$ there is a holomorphic one-parameter family of Dolbeault operators $\bar\partial_t=\bar\partial_0+tA$, $t\in\mathbb C$, such that for $t\ne0$ the holomorphic bundle defined by $\bar\partial_t$ is $\mu_P$-stable and prioritary, but not ample. As $t\to0$, these operators converge to $\bar\partial_0$ in the $C^m$ norm for every integer $m\geq0$.
\end{proposition}
\begin{proof}
Choose a nonzero extension class and a smooth Dolbeault representative $\eta\in\Omega^{0,1}(X,L_2^*\otimes L_1)$ of it, so that $\bar\partial\eta=0$. On the smooth direct sum, set
\[
 A=\begin{pmatrix}0&\eta\\0&0\end{pmatrix}.
\]
The covariant Dolbeault derivative of $A$ is zero and $A\wedge A=0$, so $\bar\partial_t^2=0$. Adding the Dolbeault operator of the base $\mathbb C$ of the parameter $t$ also gives an integrable operator on $X\times\mathbb C$, because the family depends holomorphically on $t$. This defines a holomorphic family of bundles whose extension class is $t[\eta]$. It is nonsplit for $t\neq0$, so the two preceding propositions apply. For a fixed smooth background connection and fixed metrics,
\[
 \|\bar\partial_t-\bar\partial_0\|_{C^m}
 =|t|\,\|A\|_{C^m}\longrightarrow0,
\]
and the norm $\|A\|_{C^m}$ is finite for every $m$ because $X$ is compact.
\end{proof}

% Stable-extension attribution: Liu--Wan 2609.37309v1, Proposition 4.4, pp. 20--21.
Thus stability and ampleness have been realized separately in the prioritary family of $E_0$: the nonsplit extensions are stable but not ample, while $E_{\mathrm{amp}}$ is ample. The deformation argument of \cref{thm:hirzebruch-stable-ample-approximation} joins the stable locus to the ample one and provides stable ample bundles arbitrarily close to $E_0$.
